
\SGARefTarget{sga3:xvi:section:2}%
\section*{2. A Representability Theorem for Quotients}
\addcontentsline{toc}{section}{2. A Representability Theorem for Quotients}

Let us ``recall'' the following result.

\medskip
\noindent\SGARefTarget{sga3:xvi:theorem:2:1}%
\textbf{Theorem 2.1.} --- \textit{Let \(S\) be a prescheme, let
\(X\) and \(Y\) be two \(S\)-preschemes, and let
\(f:X\to Y\) be an \(S\)-morphism. Suppose that one is in one of the
following two cases:}

\begin{enumerate}
\renewcommand{\labelenumi}{\textit{\alph{enumi})}}
\item \textit{The morphism \(f\) is locally of finite presentation.}

\item \textit{The prescheme \(S\) is locally noetherian and \(X\) is
locally of finite type over \(S\).}
\end{enumerate}

\noindent\textit{Then the following conditions are equivalent:}

\begin{enumerate}
\renewcommand{\labelenumi}{\textit{\roman{enumi})}}
\item \textit{There exist an \(S\)-prescheme \(X'\) and a
factorization of \(f\)}
\[
  f:X\xrightarrow{\,f'\,}X'\xrightarrow{\,f''\,}Y,
\]
\textit{where \(f'\) is a faithfully flat \(S\)-morphism, locally
of finite presentation, and \(f''\) is a monomorphism.}

\item \textit{The first projection}
\[
  p_1:X\times_Y X\longrightarrow X
\]
\textit{is a flat morphism.}
\end{enumerate}

\noindent\textit{Moreover, if the preceding conditions hold, then
\((X',f')\) is a quotient of \(X\), for the fpqc topology, by the
equivalence relation defined by \(f\). Consequently the
factorization \(f=f''\circ f'\) in i) is unique up to isomorphism.}

The proof of this delicate theorem will be given in EGA V. One may
also consult the lecture by J.-P. Murre, \textit{Séminaire
Bourbaki}, May 1965, no.~294, theorem~2, corollary~2, which treats
the case where \(Y\) is locally noetherian and \(X\) is of finite
type over \(Y\). We shall see that one can reduce to that case.

Let us first make a few remarks.

\medskip
\noindent\textit{a)} The implication i) \(\Rightarrow\) ii) is
trivial. Indeed, the first projection
\[
  p'_1:X\times_{X'}X\longrightarrow X
\]
factors through \(X\times_Y X\):
\[
  p'_1:X\times_{X'}X\xrightarrow{\,u\,}
  X\times_Y X\xrightarrow{\,p_1\,}X.
\]
The morphism \(u\) is an isomorphism because \(f''\) is a
monomorphism, and \(p'_1\) is flat because \(f'\) is flat; therefore
\(p_1\) is flat.

\medskip
\noindent\textit{b)} The assertions of \SGARefLink{sga3:xvi:theorem:2:1}{2.1} are local on \(Y\)
(hence local on \(S\)); they are also local on \(X\), as follows
readily from the fact that a flat morphism locally of finite
presentation is open (EGA IV, 11.3.1).

\medskip
\noindent\textit{c)} Under the hypotheses of \SGARefLink{sga3:xvi:theorem:2:1}{2.1} a), the foregoing
reduces us to the case where \(X\) and \(Y\) are affine and \(f\) is
of finite presentation. Replacing \(S\) by \(Y\), we may suppose
that \(X\) and \(Y\) are of finite presentation over \(S\). We then
reduce to the noetherian case for \(S\) by EGA IV, 11.2.6.

\medskip
\noindent\textit{d)} Under the hypotheses of \SGARefLink{sga3:xvi:theorem:2:1}{2.1} b), we may suppose
that \(S\), \(X\), and \(Y\) are affine, that \(S\) is noetherian,
and that \(X\) is of finite type over \(S\). Regard \(Y\) as a
filtered projective limit of affine schemes \(Y_i\) of finite type
over \(S\). The schemes \(X\times_{Y_i}X\) form a decreasing
filtered family of closed subschemes of \(X\times_S X\), whose
projective limit is \(X\times_Y X\). Since \(X\times_S X\) is
noetherian, one has
\[
  X\times_{Y_i}X=X\times_Y X
\]
for all sufficiently large \(i\). Thus
\[
  f_i:X\xrightarrow{\,f\,}Y\longrightarrow Y_i
\]
satisfies the hypotheses of \SGARefLink{sga3:xvi:theorem:2:1}{2.1} ii) if \(f\) does. Since the
equivalence relation on \(X\) defined by \(f\) coincides with the
one defined by \(f_i\), it is clear that it suffices to prove
ii) \(\Rightarrow\) i) for \(f_i\). This reduces us to the case
where \(Y\) is of finite type over \(S\).

\subsection*{Application to Group Preschemes}
\addcontentsline{toc}{subsection}{Application to Group Preschemes}

\noindent\SGARefTarget{sga3:xvi:theorem:2:2}%
\textbf{Theorem 2.2.} --- \textit{Let \(S\) be a prescheme and let
\(G\) be an \(S\)-group prescheme locally of finite presentation
over \(S\), acting on an \(S\)-prescheme \(X\). Let
\(\xi\in X(S)\) be a section of \(X\) over \(S\), such that the
stabilizer \(H\) of \(\xi\) in \(G\) is an \(S\)-group
subprescheme of \(G\), flat over \(S\). If \(X\) is locally of
finite type over \(S\), or if \(S\) is locally noetherian, then the
homogeneous quotient space \(G/H\) is representable by an
\(S\)-prescheme locally of finite presentation over \(S\), and the
\(S\)-morphism}
\[
  f:G\longrightarrow X,\qquad g\longmapsto g\cdot\xi
\]
\textit{factors as}
\begin{center}
  \begin{tikzcd}[column sep=5.6em,row sep=2.6em]
  G \arrow[d,"p"'] \arrow[dr,"f"]
    & {} \\
  G/H \arrow[r,"i"]
    & X.
\end{tikzcd}

\end{center}
\textit{where \(p\) is the canonical projection, a faithfully flat
morphism locally of finite presentation, and \(i\) is a
monomorphism. (For definiteness, we have supposed that \(G\) acts
on the left on \(X\).)}

\medskip
\noindent\textit{Proof.}
The morphism \(f\) makes \(G\) into an \(X\)-prescheme. By the
definition of the stabilizer of \(\xi\), the morphism
\[
  G\times_S H\longrightarrow G\times_X G,\qquad
  (g,h)\longmapsto(g,gh)
\]
is an isomorphism. Since \(H\) is flat over \(S\),
\(G\times_S H\) is flat over \(G\); hence the first projection
\[
  p_1:G\times_X G\longrightarrow G
\]
is flat. Moreover, if \(X\) is locally of finite type over \(S\),
then \(f\) is locally of finite presentation (EGA IV, 1.4.3 v));
otherwise \(S\) is locally noetherian by hypothesis. It therefore
suffices to apply \SGARefLink{sga3:xvi:theorem:2:1}{2.1} to \(f\).

It remains only to show that \(G/H\) is locally of finite
presentation over \(S\), and this follows immediately from
\SGARefLink{sga3:v:proposition:9:1}{Exposé V, 9.1}.
\hfill\(\square\)

\medskip
\noindent\SGARefTarget{sga3:xvi:corollary:2:3}%
\textbf{Corollary 2.3.} --- \textit{Let \(S\) be a prescheme, let
\(G\) and \(H\) be two \(S\)-group preschemes, and let
\(u:G\to H\) be an \(S\)-homomorphism. Suppose that \(G\) is locally
of finite presentation over \(S\), and that either \(H\) is locally
of finite type over \(S\), or \(S\) is locally noetherian. If
\(K=\Ker(u)\) is flat over \(S\), then the quotient group \(G/K\)
is representable by an \(S\)-group prescheme locally of finite
presentation over \(S\), and \(u\) factors as}
\begin{center}
  \begin{tikzcd}[column sep=5.8em,row sep=4.0em]
  G \arrow[rr,"u"] \arrow[dr,"p"']
    & {}
    & H \\
  {}
    & G/K \arrow[ur,"i"']
    & {}
\end{tikzcd}

\end{center}
\textit{where \(p\) is the canonical projection and \(i\) is a
monomorphism.}

\medskip
\noindent\textit{Proof.}
Apply \SGARefLink{sga3:xvi:theorem:2:2}{2.2} with \(X=H\) and with \(\xi\) the unit section of \(H\).
\hfill\(\square\)

\medskip
\noindent\SGARefTarget{sga3:xvi:corollary:2:4}%
\textbf{Corollary 2.4.} --- \textit{Let \(S\) be a prescheme, let
\(G\) be an \(S\)-group prescheme of finite presentation over
\(S\), let \(H\) be an \(S\)-group prescheme smooth over \(S\) with
connected fibers (and hence of finite presentation over \(S\), by
Exposé VI\textsubscript{B}, \SGARefLink{cor:sga3-VIB-5.5}{5.5}), and let \(i:H\to G\) be a
monomorphism of \(S\)-groups, so that \(H\) is a subgroup of \(G\).
Suppose that}
\[
  N=\Norm_G(H),
\]
\textit{which is representable by a closed group subprescheme of
\(G\), of finite presentation over \(S\) (\SGARefLink{sga3:xi:corollary:6:11}{Exposé XI, 6.11}), is flat
over \(S\). Then \(G/N\) is representable by an \(S\)-prescheme of
finite presentation over \(S\) and quasi-projective over \(S\).}

All the assertions to be proved, except the last one, are local on
\(S\). To establish them, we may therefore suppose that \(S\) is
quasi-compact and that the relative dimension of \(H\) over \(S\)
is constant, equal to \(r\). Proceed as in Exposé XV, \S5. For every
integer \(n>0\), let \(G^{(n)}\) (respectively \(H^{(n)}\)) be the
\(n\)-th normal invariant of the unit section of \(G\)
(respectively of \(H\)) (EGA IV, 16). The sheaf of
\(\mathcal O_S\)-modules \(H^{(n)}\) is a quotient of \(G^{(n)}\).
Since \(H\) is smooth over \(S\) of dimension \(r\), \(H^{(n)}\)
canonically defines an element \(\xi_n\) of
\[
  X_n=\operatorname{Grass}_{\varphi(n,r)}G^{(n)}
\]
(EGA I, second edition, \S9), for a suitable integer
\(\varphi(n,r)\). On the other hand, \(G\) acts naturally on
\(G^{(n)}\), and hence also on \(X_n\), through the representation
\[
  G\longrightarrow\operatorname{Aut}_{S\text{-gr}}(G),
  \qquad g\longmapsto\operatorname{int}(g).
\]
Since \(S\) is quasi-compact, there exists an integer \(m\) such
that, for \(n\geq m\),
\[
  N=\Norm_G(H)=\Norm_G H^{(n)}
  \qquad\text{(Exposé XI, 6.11 b)).}
\]
In other words, \(N\) is the stabilizer of \(\xi_n\) for sufficiently
large \(n\). The representability of \(G/N\) therefore follows from
\SGARefLink{sga3:xvi:theorem:2:2}{2.2}. The fact that \(G/N\) is of finite presentation over \(S\) is
a consequence of \SGARefLink{sga3:v:proposition:9:1}{Exposé V, 9.1}. It remains to prove that \(G/N\) is
quasi-projective over \(S\). For this, we shall exhibit a canonical
invertible sheaf on \(G/N\) that is ample relative to \(S\).
Consider the functor
\[
  \mathscr F:(\mathbf{Sch}/S)^\circ\longrightarrow\mathbf{Ens}
\]
such that, for every \(S\)-prescheme \(T\),
\[
  \mathscr F(T)=
  \left\{
  \begin{array}{c}
  \text{\(T\)-subgroups \(H\) of \(G_T\), representable and smooth
  over \(T\),}\\
  \text{with connected fibers}
  \end{array}
  \right\}.
\]

Such a group \(H\) is of finite presentation over \(S\)
(Exposé VI\textsubscript{B}, 5.5), and
\(H\to G_T\) is a quasi-affine morphism (EGA IV, 8.11.2).
By effective descent for quasi-affine morphisms (SGA 1, VIII, 7.9),
it follows that \(\mathscr F\) is a sheaf for the fpqc topology.
Since \(G/N\) is the sheaf associated with a subfunctor of
\(\mathscr F\), there is a canonical monomorphism
\[
  G/N\longrightarrow\mathscr F.
\]
Consequently there exists a subgroup \(H'\) of
\(G\times_S(G/N)\), representable and smooth over \(G/N\), with
connected fibers, that is ``universal'' for the functor \(G/N\).
I claim that the invertible sheaf
\[
  \mathscr L=\bigl(\det(\operatorname{Lie}H')\bigr)^{-1}
\]
is ample relative to \(S\). In this form the assertion is local on
\(S\), and its proof is analogous to the one given in \SGARefLink{sga3:xv:remark:5:8}{Exposé XV,
5.8}.

The following corollary was announced in \SGARefLink{sga3:xiv:4.8.1}{Exposé XIV, 4.8 bis}.

\medskip
\noindent\SGARefTarget{sga3:xvi:corollary:2:5}%
\textbf{Corollary 2.5.} --- \textit{Let \(S\) be a prescheme, let
\(G\) be an \(S\)-group prescheme smooth and of finite presentation
over \(S\), with connected fibers, and let \(P\) be a parabolic
subgroup of \(G\) (\SGARefLink{sga3:xiv:paragraph:4-8-bis}{Exposé XIV, 4.8 bis}). Then \(G/P\) is
representable by an \(S\)-prescheme smooth and projective over
\(S\).}

It remains only to show (\textit{loc.\ cit.}) that \(G/P\) is
representable by an \(S\)-prescheme that is quasi-projective and of
finite presentation. This follows from \SGARefLink{sga3:xvi:corollary:2:4}{2.4}, since
\[
  P=\Norm_G(P)
\]
(\textit{loc.\ cit.}) and is therefore flat over \(S\).
\hfill\(\square\)
